\hypertarget{harmonic-returns-harmonic-scale-return-laws-for-the-kicked-pell-cat-map}{%
\section{Harmonic-scale return laws for the kicked Pell cat map}\label{harmonic-returns-harmonic-scale-return-laws-for-the-kicked-pell-cat-map}}

\textbf{Status: project proof.} This is a global spectral theorem for the actual
return propagator. It concerns a vanishing quantum kick and times of order
the arithmetic period, not Heisenberg-scale correlations or BGS. The endpoint
probability laws are known; priority for their arithmetic realization is not
established. All asymptotic statements below keep the return multiplicity
and kick parameter fixed.

\hypertarget{harmonic-returns-statement}{%
\subsection{1. Statement}\label{harmonic-returns-statement}}

Use the quantization and physical harmonics of
the preceding harmonic observable theorem:
\[
 N+F\sqrt2=(1+\sqrt2)^K,\qquad D=(N-1)/2,\qquad U_0^K=I,
 \qquad B_a(x)=2\cos(2\pi ax/N).
\]
Here \(K\) runs through odd primes. Let \(a\) be uniform on
\(\mathcal A_K=\{1\le a\le L_K:\gcd(a,N)=1\}\), where
\(1\le L_K\le\lfloor\sqrt N\rfloor\) and
\(\log L_K/\log N\to\beta\in[0,1/2]\). Put
\[
 \epsilon=\frac{\tau}{\sqrt{2K}},\qquad
 U_{\epsilon,a}=e^{i\epsilon B_a}U_0,\qquad t=\tau^2,\qquad
 w_\beta(x,y)=\mathbf1\{\operatorname{dist}(x,y)<(1-2\beta)/2\}.
                                                               \tag{R1}
\]
The normalized trace \(\operatorname{tr}_-\) is on the full odd parity space.

For \(m\ge1\), let \(\mathcal C_m=\{(r,s):0\le r<s<m\}\).
Given \(\ell\) ordered labels and types
\(c_i=(r_i,s_i)\in\mathcal C_m\), arrange their \(2\ell\) endpoints
\((r_i,i),(s_i,i)\) in lexicographic order (cycle first, label second).
Join the two endpoints of each label. Write \(G(c)\) for the crossing graph
of these chords, retaining the label \(i\) on its vertex. Set
\[
 I_c(\beta)=\int_{0<x_1<\cdots<x_\ell<1}
                  \prod_{\{i,j\}\in E(G(c))}w_\beta(x_i,x_j)\,dx,
\]
\[
 F_{m,\beta}(t)=e^{-mt/2}\sum_{\ell=0}^\infty(-t)^\ell
                  \sum_{c\in\mathcal C_m^\ell} I_c(\beta).       \tag{R2}
\]
The empty integral is one. For \(m=1\) there are no nonempty type lists.
The series is absolutely convergent, since its \(\ell\)th absolute term is
at most \((\binom m2 t)^\ell/\ell!\).

\begin{theorem}\label{pass1:pell_harmonic_returns-1} For every fixed real \(\tau\) and positive integer \(m\),
\[
 \left\|\operatorname{tr}_- U_{\epsilon,a}^{mK}
                 -F_{m,\beta}(\tau^2)\right\|_{L^2(\mathcal A_K)}
       \longrightarrow0.                                    \tag{R3}
\]
There is a unique probability measure \(\mu_{\beta,t}\) on the unit circle
with positive Fourier moments \(F_{m,\beta}(t)\). The empirical eigenvalue
measures of \(U_{\epsilon,a}^K\) converge weakly to it in probability, and
their averages converge weakly. In particular
\[
 F_{1,\beta}(t)=e^{-t/2},\qquad
 F_{m,0}(t)=e^{-m^2t/2},\qquad
 F_{m,1/2}(t)=\frac{e^{-mt/2}}m L_{m-1}^{(1)}(mt).              \tag{R4}
\]
Thus the endpoints are the wrapped normal law and the spectral law of free
unitary Brownian motion at time \(t\), respectively. The latter identification
uses the known moment formula; the calculation of these moments from the
quantum map is proved below. This is not the exponential of a semicircular
random variable.
\end{theorem}

There is also a closed intermediate formula:
\[
 F_{2,\beta}(t)=e^{-t}\left(1-t e^{-(1-2\beta)t/2}\right),
                      \qquad 1/6\le\beta\le1/2.               \tag{R5}
\]
The family \(\mu_{\beta,t}\) is weakly continuous in \((\beta,t)\).

\hypertarget{harmonic-returns-exact-expansion-and-uniform-truncation}{%
\subsection{2. Exact expansion and uniform truncation}\label{harmonic-returns-exact-expansion-and-uniform-truncation}}

Conjugation and dilation as in the harmonic theorem turn the physical kick
into \(W(q,1)+W(q,1)^*\) with \(q=a^2q_0\), \(q_0=-F\). The square of \(a\) is
essential. The unperturbed operator becomes \(S_h^{-1}\), \(h=3+2F^{-1}\).
Choose orbit vectors \(v_j=(qh^j,h^{-j})\), \(0\le j<K\); these have precisely
the order in
\[
 U_{\epsilon,a}^{mK}
   =\left[\prod_{j=0}^{K-1}
      \exp\{i\epsilon(W(v_j)+W(-v_j))\}\right]^m.               \tag{R6}
\]
All products are read left to right with increasing \(j\).

The absolutely convergent Fourier--Bessel identity is
\[
 e^{i\epsilon(W+W^*)}=\sum_{n\in\mathbb Z}i^nJ_n(2\epsilon)W^n,
 \qquad |J_n(2\epsilon)|\le I_n(2|\epsilon|).                   \tag{R7}
\]
It follows directly by multiplying the power series for
\(e^{i\epsilon W}\) and \(e^{i\epsilon W^{-1}}\); thus it applies equally to
finite-order Weyl unitaries. Let \(n_{r,j}\) be the coefficient chosen in
cycle \(r\), put \(R=\sum_{r,j}|n_{r,j}|\), and
\(c_j=\sum_r n_{r,j}\). A word is a scalar phase times
\(W(\sum_j c_jv_j)\).

Individual Weyl operators do not preserve parity. Applying the full-space
trace with the odd projector gives, nonetheless, the exact linear identity
\[
 D^{-1}\operatorname{Tr}\left(\frac{I-\mathcal R}{2}W(v)\right)
             =\frac{N\mathbf1_{\{v=0\}}-1}{N-1}.               \tag{R8}
\]
Indeed \(\operatorname{Tr}W(v)=N\mathbf1_{\{v=0\}}\) and
\(\operatorname{Tr}(\mathcal R W(v))=1\) for the centered convention.
Hence replacing every normalized word trace by its central indicator
times its phase has total error at most
\[
 \frac{e^{2mK\delta}}{N-1}=o(1),\qquad \delta=|\epsilon|.       \tag{R9}
\]
Here \(\sum_n I_n(2\delta)=e^{2\delta}\).

Let \(R_0=\lfloor K^{2/3}\rfloor\). For \(z\ge1\), the coefficient series gives
\[
 I_n(2\delta)\le e^{\delta^2}\frac{\delta^{|n|}}{|n|!},
 \qquad
 \sum_n I_n(2\delta)z^{|n|}
          \le e^{\delta^2}(2e^{\delta z}-1)
          \le e^{\delta^2+2\delta z}.
\]
Thus the absolute weight of words with \(R>R_0\) is at most
\[
 z^{-R_0}\exp(mK\delta^2+2mK\delta z).
\]
For \(\tau\ne0\) take \(z=R_0/(2mK\delta)>1\), eventually. This bound is
\[
 \exp\{mK\delta^2+R_0-R_0\log(R_0/(2mK\delta))\}
 =\exp\{-(1/6+o(1))K^{2/3}\log K\}=o(1).                      \tag{R10}
\]
For \(\tau=0\) the expansion has only its constant term.

\hypertarget{harmonic-returns-all-order-suppression-of-nonzero-aggregate-relations}{%
\subsection{3. All-order suppression of nonzero aggregate relations}\label{harmonic-returns-all-order-suppression-of-nonzero-aggregate-relations}}

\begin{lemma}\label{pass1:pell_harmonic_returns-2} Uniformly in unit \(a\), the absolute weight of central words
with \(R\le R_0\) and \(c\ne0\) is \(O_{m,\tau}(K^{-1})\).
\end{lemma}

\textbf{Proof.} Centrality and the unit property of \(q\) give both congruences
\(\sum c_jh^j=\sum c_jh^{-j}=0\). Cancel zero coefficients and cut the support
at a largest cyclic gap. Its size is at least \(\lceil K/R\rceil\).
The rotated polynomial \(P\) has coefficient norm at most \(R\) and degree
at most \(K-\lceil K/R\rceil\). Dividing by the monic polynomial
\(z^2-6z+1\), the two congruences and the unit \(h-h^{-1}\) show that both
remainder coefficients are divisible by \(N\). A nonzero quadratic norm
would therefore obey
\[
 N^2\le|\operatorname{Norm}P(\lambda)|
          \le R^2\lambda^{K-\lceil K/R\rceil},
                  \qquad \lambda=3+2\sqrt2.                   \tag{R11}
\]
This is impossible uniformly for \(R\le R_0\): since
\(N^2=\lambda^K(1-\lambda^{-K})^2/4\), the ratio is at most
\(4R_0^2\lambda^{-\lceil K/R_0\rceil}/(1-\lambda^{-K})^2=o(1)\).
Consequently \(P\) is an integer multiple of \(z^2-6z+1\).

Work now in \(\mathbb Z[z]/(z^K-1)\). Rotation back is multiplication by
a power of \(z\), so the original coefficient vector satisfies
\[
 c=Mq',\qquad M=S^2-6S+I,\qquad q'\in\mathbb Z^K,              \tag{R12}
\]
where \(S\) is cyclic shift. The triangle inequality gives
\[
 \|Mq'\|_1\ge6\|Sq'\|_1-\|S^2q'+q'\|_1\ge4\|q'\|_1.
                                                               \tag{R13}
\]
In particular \(M\) is injective; no counting multiplicities of preimages
is needed.

For a prescribed \(c\), summing absolute weights over its cycle assignments
is bounded, by modified-Bessel convolution, by
\[
 \prod_{j=0}^{K-1} I_{c_j}(2m\delta)
       \le e^{Kb^2}b^{\|c\|_1},\qquad b=m\delta<1.
\]
We may enlarge the truncated sum to all assignments with \(c\) in the
integer image of \(M\). Equations (R12)--(R13) imply the bound
\[
 e^{Kb^2}\sum_{q'\ne0}b^{4\|q'\|_1}
  =e^{Kb^2}\left[\left(\frac{1+b^4}{1-b^4}\right)^K-1\right]
  =O_{m,\tau}(K^{-1}).                                      \tag{R14}
\]
This proves the lemma. In contrast to fixed-order relation counting,
(R14) controls the entire truncated exponential at once.

\hypertarget{harmonic-returns-balanced-sites-and-the-limiting-series}{%
\subsection{4. Balanced sites and the limiting series}\label{harmonic-returns-balanced-sites-and-the-limiting-series}}

After (R9)--(R14), only coefficientwise-balanced arrays \(c=0\) remain.
At a site \(j\) the total absolute coefficient majorant is \(I_0(2m\delta)\).
The inactive array has majorant \(I_0(2\delta)^m\); therefore the active
majorant is
\[
 A=I_0(2m\delta)-I_0(2\delta)^m
             =m(m-1)\delta^2+O_m(\delta^4).
\]
In particular, the total mass with \(\ell\) active sites is at most
\[
 \binom K\ell A^\ell I_0(2\delta)^{m(K-\ell)}
       \le e^{O_{m,\tau}(1)}\frac{C_{m,\tau}^{\,\ell}}{\ell!}.  \tag{R15}
\]
The bound is summable uniformly in sufficiently large \(K\).

An active balanced site has even total order. Its order-two arrays have
one \(+1\) and one \(-1\) in two different cycles. The majorant of all
order-at-least-four balanced arrays is \(O_m(\delta^4)\), by its convergent
coefficient series. Marking one such site and summing the remaining sites
gives an error
\(O_m(K\delta^4)I_0(2m\delta)^{K-1}=O_{m,\tau}(K^{-1})\).
Every remaining active site is consequently a chord with a type in
\(\mathcal C_m\) and one of two orientations. Its exact signed Bessel
coefficient is \(-J_1(2\epsilon)^2\), because
\(i^{-1}J_{-1}(2\epsilon)=iJ_1(2\epsilon)\).
For \(\ell\) such sites the common coefficient is
\[
 (-J_1(2\epsilon)^2)^\ell J_0(2\epsilon)^{mK-2\ell}.           \tag{R16}
\]

Fix \(\ell\), its type list and orientations. The Weyl phase is a sum of
crossing contributions. Indexing chords in first-endpoint order, as in
(H9) of the harmonic theorem, (R6) uses the inverse orbit and reverses
the phase sign. If chords are instead indexed by time label, each edge
also records which chord has the earlier first endpoint. These signs
do not change the limiting argument. The phase lemma of that theorem
extends to the ordered grid \(j_1<\cdots<j_\ell\) as follows. Remove
neighborhoods of coincident labels, circle-distance boundaries and
equal crossing distances, of normalized size \(O_\ell(\eta+K^{-1})\).
On the remaining configurations, the unique largest crossing distance
dominates the signed Pell sum. Its harmonic average tends to one below
the threshold \((1-2\beta)/2\) and to zero above it, by the quadratic-sum
estimate (H11)--(H13). The same argument on two copies, additionally
excluding ties between edges in distinct copies, gives variance tending
to zero. Restriction to the simplex merely inserts a bounded indicator
whose boundary has measure zero. Thus
\[
 K^{-\ell}\sum_{0\le j_1<\cdots<j_\ell<K}
       \text{phase}(a;j,c,\sigma)
       \longrightarrow I_c(\beta)\quad\hbox{in }L^2(\mathcal A_K).
                                                               \tag{R17}
\]
For edgeless diagrams the phase is identically one, and
\(K^{-\ell}\binom K\ell\to1/\ell!\), including \(\ell=0\).
Conditioning on units has vanishing
cost by (H8); all these normalized phase sums are bounded.

Now \(KJ_1(2\epsilon)^2\to t/2\) and
\(J_0(2\epsilon)^{mK-2\ell}\to e^{-mt/2}\) at fixed \(\ell\).
The \(2^\ell\) orientations cancel the factors \(2^{-\ell}\).
Equations (R15)--(R17), with a fixed active-site cutoff followed by
its removal, prove (R2)--(R3). The deterministic absolute errors are
uniform in \(a\); the fixed-diagram phase limits supply the \(L^2\) convergence.

The Fourier moments of averaged empirical measures converge to (R2).
Compactness of the circle supplies a subsequential probability limit;
all its Fourier moments are specified, including the negative ones by
complex conjugation. Trigonometric-polynomial approximation gives
uniqueness and convergence of the full sequence. The same approximation
and (R3) give weak convergence in probability. Dominated convergence in
the simplex integrals and in (R2) gives continuity in \(\beta,t\), uniformly
on compact \(t\) intervals at each fixed \(m\).

\hypertarget{harmonic-returns-endpoints-a-direct-noncrossing-enumeration}{%
\subsection{5. Endpoints: a direct noncrossing enumeration}\label{harmonic-returns-endpoints-a-direct-noncrossing-enumeration}}

At \(\beta=0\), all \(I_c=1/\ell!\). Summation in (R2) gives
\(e^{-mt/2}e^{-\binom m2 t}=e^{-m^2t/2}\).
At \(\beta=1/2\), only noncrossing type lists survive and each contributes
\(1/\ell!\). Let \(A_{m,\ell}\) count these lists; \(A_{m,0}=1\).

For \(\ell\ge1\), delete the chord with the largest label, of type \((r,s)\).
Its endpoints split the remaining circular order into complete sets of
\(s-r\) and \(m-s+r\) cycle rows. Noncrossing forces every other chord to stay
in one of the two arcs. If \(p\) of the remaining labels lie in the first
arc, their choice gives \(\binom{\ell-1}{p}\); both relative orders retain
the original definition. Conversely any two such noncrossing lists give
exactly one list with the marked chord restored. Hence
\[
 A_{m,\ell}=\frac m2\sum_{j=1}^{m-1}\sum_{p=0}^{\ell-1}
       \binom{\ell-1}{p}A_{j,p}A_{m-j,\ell-1-p}.               \tag{R18}
\]
The coefficient is initially \(m-j\), the number of \((r,s)\) with difference
\(j\); symmetry under \(j\leftrightarrow m-j\) replaces it by \(m/2\).

Set \(H_m(t)=\sum_\ell A_{m,\ell}(-t)^\ell/\ell!\).
Equation (R18) gives \(H'_m=-(m/2)\sum_{j=1}^{m-1}H_jH_{m-j}\),
\(H_m(0)=1\). As formal series in \(z\), \(H=\sum_{m\ge1}H_mz^m\) obeys
\(\partial_tH=-zH\partial_zH\), \(H(0,z)=z/(1-z)\), or equivalently
\[
 H=z(1+H)e^{-tH}.
\]
Formal Lagrange inversion then yields
\[
 H_m(t)=\frac1m[u^{m-1}](1+u)^m e^{-mtu}
   =\frac1m\sum_{\ell=0}^{m-1}\binom m{\ell+1}
                           \frac{(-mt)^\ell}{\ell!}.
                                                               \tag{R19}
\]
The last polynomial is \(L_{m-1}^{(1)}(mt)/m\). This proves (R4),
including \(A_{m,\ell}=m^{\ell-1}\binom m{\ell+1}\) for
\(0\le\ell<m\) and zero for \(\ell\ge m\).
For the conventional identification with free unitary Brownian motion,
see N. Demni, \emph{Lagrange inversion formula, Laguerre polynomials and the
free unitary Brownian motion}, equation (1),
\href{https://arxiv.org/pdf/1606.03755v3}{arXiv:1606.03755v3}.
No stochastic-process convergence is inferred from the marginal law.

\hypertarget{harmonic-returns-a-closed-second-return-formula}{%
\subsection{6. A closed second-return formula}\label{harmonic-returns-a-closed-second-return-formula}}

For \(m=2\) every type is \((0,1)\) and all distinct chords cross.
Writing \(b=(1-2\beta)/2\), we obtain
\[
 F_{2,\beta}(t)=e^{-t}\sum_{\ell=0}^\infty
                  \frac{(-t)^\ell}{\ell!}\,t(K_\ell,w_\beta).
                                                               \tag{R20}
\]
When \(0<b<1/3\), a finite set of circle points at pairwise distance less
than \(b\) must lie in an arc of length \(b\). To see this, put one point at
zero. All others lie within clockwise distance \(b\) or counterclockwise
distance \(b\). If the maximal distances in these two directions are \(u,v\),
then \(u+v\le2b<1-b\); pairwise closeness therefore forces \(u+v<b\).
The covering arc has a unique first sample point, except on a null set.
For \(\ell\) independent uniform points, each candidate first point covers
all others with probability \(b^{\ell-1}\), and these events are disjoint.
Consequently \(t(K_\ell,w_\beta)=\ell b^{\ell-1}\) for \(\ell\ge1\).
The formulas extend to \(b=1/3\) and \(b=0\) by continuity, interpreting
the \(\ell=1\) value as one. Summing (R20) proves (R5).

\hypertarget{harmonic-returns-what-this-removes-and-what-remains-open}{%
\subsection{7. What this removes, and what remains open}\label{harmonic-returns-what-this-removes-and-what-remains-open}}

The earlier weighted-observable moment limit alone did not justify
time-ordered exponentiation. Equations (R10)--(R15) provide the missing
uniform summability for this arithmetic model. They also avoid extending
the small-harmonic commuting normal form beyond its valid phase regime.

The new observable is the actual unitary \(U_{\epsilon,a}^K\). Nevertheless
\(K\asymp\log N\), whereas the relevant Heisenberg scale is of order \(D\).
Here \(m\) is fixed, and constants and phase-cutoff complexities depend on
\(m\). Neither this proof nor its noncrossing endpoint controls
\(m\asymp D/K\), joint microscopic correlations, or fixed-strength classical
perturbations. In particular the result supplies no positive BGS theorem.
Lean formalization and an external referee assessment remain outstanding.
